% Лабораторная работа № 7. Компиляция: pdflatex report_lab7.tex дважды.
\documentclass[a4paper,14pt]{extarticle}

\usepackage[utf8]{inputenc}
\usepackage[T2A]{fontenc}
\usepackage[russian]{babel}
\usepackage{geometry}
\usepackage{setspace}
\usepackage{indentfirst}
\usepackage{titlesec}
\usepackage{tocloft}
\usepackage{enumitem}
\usepackage{booktabs}
\usepackage{array}
\usepackage{longtable}
\usepackage{graphicx}
\usepackage{fancyvrb}
\usepackage[hidelinks,unicode]{hyperref}

\geometry{left=30mm,right=15mm,top=20mm,bottom=20mm}
\onehalfspacing
\setlength{\parindent}{12.5mm}
\setlength{\parskip}{0pt}
\setlength{\emergencystretch}{3em}
\renewcommand{\contentsname}{СОДЕРЖАНИЕ}
\renewcommand{\refname}{СПИСОК ИСПОЛЬЗОВАННЫХ ИСТОЧНИКОВ}
\titleformat{\section}{\centering\bfseries\fontsize{14pt}{16pt}\selectfont}{\thesection}{1em}{}
\titleformat{\subsection}{\bfseries\fontsize{14pt}{16pt}\selectfont}{\thesubsection}{1em}{}
\titlespacing*{\section}{0pt}{12pt}{12pt}
\titlespacing*{\subsection}{\parindent}{12pt}{6pt}
\setlist{nosep}
\setcounter{tocdepth}{2}
\hypersetup{pdftitle={Отчёт по лабораторной работе № 7},pdfauthor={Студент группы P3113}}

\begin{document}
\begin{titlepage}
\begin{center}
Федеральное государственное автономное образовательное учреждение\\
высшего образования\\
\textbf{«Национальный исследовательский университет ИТМО»}\\[10mm]

Факультет программной инженерии и компьютерной техники\\
Дисциплина «Программирование»\\[34mm]

{\large Лабораторная работа № 7}\\[4mm]
Многопользовательское клиент-серверное приложение\\
Вариант \textbf{9585}
\end{center}

\vfill
\begin{flushright}
Выполнил:\\
Студент группы P3113
Ушаков Роман Павлович\\[5mm]
Проверил:\\
Гиря Максим Дмитриевич

\end{flushright}

\vfill
\begin{center}
г. Санкт-Петербург\\ 2026
\end{center}
\end{titlepage}

\setcounter{page}{2}
\tableofcontents
\newpage

\section{Задание}

Переработать клиент-серверное приложение из лабораторной работы №6. Сервер должен хранить коллекцию в базе данных PostgreSQL, поддерживать регистрацию и авторизацию пользователей, а также выполнять изменяющие команды только для объектов их владельцев. Необходимо обеспечить потокобезопасную работу сервера и журналирование его действий.

\section{Цель работы}

Освоить работу с реляционными базами данных через JDBC, реализацию аутентификации и разграничения доступа в клиент-серверном приложении, а также многопоточную обработку запросов.

\section{Используемые средства}

Язык реализации --- Java~17. Сборка выполняется Apache Maven. Для постоянного хранения коллекции используется PostgreSQL и JDBC-драйвер. Передача сериализованных объектов выполняется по UDP с помощью \texttt{DatagramChannel}. Пароли пользователей хранятся в виде SHA-224-хэшей. Журналирование реализовано библиотекой Log4j2.

\section{Диаграмма классов UML}

UML-диаграмма основных классов вынесена в отдельный файл \href{uml_class_diagram.pdf}{\texttt{uml\_class\_diagram.pdf}}. Её исходник в TikZ находится в файле \texttt{uml\_class\_diagram.tex} рядом с данным документом.

\section{Исходный код программы}

Исходный код и SQL-скрипты находятся в каталоге лабораторной работы в репозитории. В каталоге \texttt{sql} расположены схема базы данных, скрипт начального заполнения и проверочные запросы.

\section{Результат работы программы}

Программа разделена на модули клиента, сервера и общих классов. При запуске сервер через JDBC загружает коллекцию \texttt{SpaceMarine} из PostgreSQL в потокобезопасные структуры данных. Команды чтения работают с коллекцией в памяти; изменяющие операции сначала выполняются в базе данных и после успешного SQL-запроса отражаются в памяти.

Клиент поддерживает команды \texttt{register} и \texttt{login}. После авторизации пользователь может просматривать всю коллекцию, а изменять и удалять --- только созданные им объекты. Идентификатор и дата создания нового объекта назначаются PostgreSQL. Сервер принимает UDP-запросы в отдельном потоке, создаёт поток для обработки каждого запроса и отправляет ответы через кэшируемый пул потоков. Ход работы и ошибки записываются Log4j2 в консоль и журнал запуска.

\end{document}
